% Tabel legenda level kematangan MLOps (Google Cloud), portrait, lintas halaman
% Non-berskor; mengikuti gaya longtable mlops_maturity_comparison.tex + rubrik Lampiran A
% Karakteristik dan komponen verbatim dari googlecloud2024mlops
{
\footnotesize
\setlength{\tabcolsep}{3pt}
\renewcommand{\arraystretch}{1.2}
\setlength{\LTpost}{0pt}
\Needspace{13\baselineskip}
\begin{longtable}{|>{\centering\arraybackslash}m{1.5cm}|>{\raggedright\arraybackslash}m{5.95cm}|>{\raggedright\arraybackslash}m{5.8cm}|}
\caption{Karakteristik dan Komponen Tiap \textit{Level} Kematangan MLOps}
\label{tab:mlops_level_legend} \\
\hline
\textbf{\textit{Level}} & \centering\arraybackslash\textbf{Karakteristik Proses} & \centering\arraybackslash\textbf{Komponen} \\
\hline
\endfirsthead
\hline
\textbf{\textit{Level}} & \centering\arraybackslash\textbf{Karakteristik Proses} & \centering\arraybackslash\textbf{Komponen} \\
\hline
\endhead
\endfoot
\hline
\endlastfoot

\textbf{\textit{Level} 0}\newline \textit{Manual process} & \textit{Manual, script-driven, and interactive process} dengan \textit{disconnection between ML and operations} serta \textit{infrequent release iterations}. Tanpa CI dan tanpa CD, \textit{deployment} hanya merujuk pada \textit{prediction service}, disertai \textit{lack of active performance monitoring}. & Tanpa komponen otomasi tambahan. Penyerahan model bersifat satu arah dari \textit{data scientist} ke tim operasi, dan \textit{deployment} terbatas sebagai \textit{prediction service}. \\
\hline
\textbf{\textit{Level} 1}\newline \textit{ML pipeline automation} & \textit{Rapid experiment}, \textit{CT of the model in production}, dan \textit{experimental-operational symmetry}. \textit{Modularized code for components and pipelines}, \textit{continuous delivery of models}, serta \textit{pipeline deployment}. & Komponen tambahan yang dituntut \textit{level} ini, yaitu \textit{data and model validation}, \textit{feature store} (opsional), \textit{metadata management}, dan \textit{ML pipeline triggers}. \\
\hline
\textbf{\textit{Level} 2}\newline \textit{CI/CD pipeline automation} & Enam tahap otomasi \textit{pipeline}, yaitu \textit{development and experimentation}, \textit{pipeline continuous integration}, \textit{pipeline continuous delivery}, \textit{automated triggering}, \textit{model continuous delivery}, dan \textit{monitoring}. & Tujuh komponen wajib, yaitu \textit{source control}, \textit{test and build services}, \textit{deployment services}, \textit{model registry}, \textit{feature store}, \textit{ML metadata store}, dan \textit{ML pipeline orchestrator}. \\
\hline
\end{longtable}
\par{\footnotesize Sumber: \textcite{googlecloud2024mlops}. Istilah teknis dipertahankan dalam bahasa asli sesuai dokumen sumber.}
}
